
\newpage
\section*{RESUMEN} % Se añade un asterisco a \section para que el título no esté numerado.
\markright{RESUMEN} % Al utilizar \section* se ha de añadir manualmente el título del apartado al encabezado.
\addcontentsline{toc}{section}{RESUMEN} % Al utilizar \section* se ha de añadir manualmente el apartado al índice (Table Of Contents, TOC)

Las interfaces cerebro-computador dedicadas al control de exoesqueletos constituyen una herramienta prometedora para la rehabilitación de personas con lesiones motoras. En este contexto, la detección precisa de la intención de marcha del usuario a partir de señales electroencefalográficas (EEG) resulta fundamental para garantizar un funcionamiento seguro y eficiente del sistema. Así pues, el desarrollo de modelos de clasificación robustos y generalizables puede ayudar a reducir la dependencia de calibraciones extensas y mejorar la experiencia del usuario, facilitando la integración de estas tecnologías en entornos clínicos y domésticos.

El objetivo de este trabajo es analizar el rendimiento de diferentes arquitecturas de redes neuronales convolucionales en la clasificación de señales EEG, así como evaluar distintas estrategias de entrenamiento, incluyendo modelos entrenados con datos del mismo día, enfoques acumulativos y técnicas de fine-tuning. Entre las arquitecturas evaluadas se incluyen \textit{EEGNet}, \textit{DeepConvNet} y \textit{ShallowConvNet}, mientras que las estrategias de entrenamiento consideradas abarcan enfoques \textit{Intra-Sesión}, \textit{Inter-Sesión}, \textit{Inter-Usuarios} y \textit{Continuo}.

Los experimentos se han realizado utilizando registros de señales EEG correspondientes a cinco usuarios y cinco sesiones por participante. Las tareas objetivo han sido la detección de la intención de caminar y detenerse. Las diferentes configuraciones se han entrenado con varias semillas aleatorias y se han evaluado mediante la precisión media y su desviación estándar, permitiendo estudiar tanto el rendimiento como la estabilidad de los modelos. De esta manera, se ha analizado la influencia de la arquitectura, la tarea, el método de entrenamiento y la variabilidad entre usuarios.

Los resultados muestran que \textit{ShallowConvNet} obtiene el mejor rendimiento global, con una precisión media aproximada del 76.6\%, seguida de \textit{EEGNet}, con un 71.4\%, y \textit{DeepConvNet}, con un 69.4\%. La tarea estática alcanza una precisión media cercana al 76.0\%, frente al 69.1\% obtenido en la tarea de movimiento. En cuanto a las estrategias de entrenamiento, el método \textit{Inter-Sesión} presenta el mejor resultado medio, aunque las diferencias entre métodos son moderadas y dependen en gran medida de la arquitectura, la tarea y el participante.

En conjunto, los resultados indican que la arquitectura seleccionada y la variabilidad entre usuarios tienen una mayor influencia sobre el rendimiento que la estrategia de entrenamiento. Este trabajo proporciona una base para el desarrollo de pruebas más exhaustivas orientadas a reducir el proceso de calibración de los sistemas de asistencia robótica y facilitar su integración en aplicaciones de rehabilitación.



\afterpage{\blankpage} % Se añade una página en blanco después del resumen.
